% No preamble. Prefix cdmom-. Polynomial jets, not a periodic flux family.
\section{Constant Cartan curvature and the actual stationary jets}
\label{cdmom-section}

We now use the diagonal-density construction to make a
constant-curvature calculation an exact test of the actual
Maxwell coefficient. All substitutions at an unbounded
affine field mean substitutions into exponentially localized
Taylor kernels at zero, through the required finite order.

\subsection{Exact local Wilson and averaging identities}
\label{cdmom-affine-section}

Choose commuting traceless Hermitian matrices \(c_\mu,B_{\mu\nu}\)
in a Cartan subalgebra of \(\mathfrak{su}(N)\), with
\(B_{\mu\mu}=0\), and put
\begin{equation}
 A_\mu(x)=c_\mu+\sum_\nu B_{\mu\nu}x_\nu,\qquad
 U_\mu(x)=e^{iA_\mu(x)}.
 \label{cdmom-affine}
\end{equation}
The plaquette in the \(\mu,\nu\) plane is the constant matrix
\[
 e^{iF_{\mu\nu}},\qquad F_{\mu\nu}=B_{\nu\mu}-B_{\mu\nu}.
\]
The full Wilson Euler derivative is zero. For each varied
bond and transverse direction, the two incident plaquettes
give opposite transported sine terms. They agree after
transport because all background links and plaquettes
commute. The cancellation occurs as a Lie-algebra-valued
identity, including against variations outside the Cartan
subalgebra. This is a local infinite-lattice identity;
an artificial finite-box boundary is not included.

For a centered primitive block, write \(x=y+r\).
Every coordinate-order root contour has transporter
\[
 \exp\{i[r\cdot A(y)+q_\Gamma(B,r)]\},
\]
where \(q_\Gamma\) is independent of \(y\).
The selected mean is endpoint equivariant and conjugation
equivariant. Its value on a tuple of diagonal matrices
is diagonal: conjugation by every diagonal group element
fixes the tuple, hence its output. Therefore the actual
root transporter has the form
\begin{equation}
 P_\ell(U;y,y+r)=e^{ir\cdot A(y)}p_\ell(B,r).
 \label{cdmom-contour}
\end{equation}
No log-arithmetic formula for the selected mean is assumed.

The \(p_\ell\) factors cancel in the ordinary path product
between neighboring blocks. Since \(B_{\mu\mu}=0\), this
product is
\[
 U_x=e^{i\ell A_\mu(y)}
       \exp\!\left\{i\ell\sum_\nu
                 (B_{\mu\nu}-B_{\nu\mu})r_\nu\right\}.
\]
Its reference center string is \(e^{i\ell A_\mu(y)}\).
The defining ordinary average takes the centered average
of the relative logarithms, and \(\sum_r r_\nu=0\).
It follows that the actual primitive and iterated maps obey
\begin{equation}
 M_\ell(U)(y,y+\ell e_\mu)=e^{i\ell A_\mu(y)},\qquad
 F_K(U)(y,y+Ke_\mu)=e^{iK A_\mu(y)},\quad K=\ell^k.
 \label{cdmom-average}
\end{equation}
Thus the output affine curvature is \(K^2F\).
These are identities of local analytic germs, hence
identities of formal curvature jets at each lattice site.
They do not require a common bounded logarithm of the
affine field on the whole infinite lattice.

Apply the actual direct-\(K\) root gauge
\(u(x)=P_K(U;y,x)\) in the block centered at \(y\).
It is the identity at centers. Endpoint covariance makes
the transformed contour equal to the identity and leaves
the composite output unchanged. Wilson stationarity
is preserved by gauge invariance. At every fixed order
in a curvature parameter, the rooted candidate has
polynomial growth in the block center. The within-block
coordinate set is finite at fixed \(K\).

\subsection{Why the formal candidate is the actual minimum jet}
\label{cdmom-jet-section}

\begin{lemma}[Polynomial extension of the stationary recursion]
\label{cdmom-polynomial-lemma}
Suppose an actual stationary equation in retained coordinates is
\[
 \mathcal A z+\mathcal B w+\mathcal N(w,z)=0,\qquad
                     \mathcal N=O((w,z)^2).
\]
Assume \(\mathcal A,\mathcal A^{-1},\mathcal B\) have
exponentially summable kernels, the inverse identities
hold on finitely supported fields, and the mixed Taylor
kernels of \(\mathcal N\) through a specified finite
order have exponential rooted-tree bounds. Then its
analytic solution jets extend uniquely to polynomial-growth
input jets. Any polynomial-growth formal solution has
those same coefficients through that order.
\end{lemma}

\begin{proof}
For a polynomial-growth sequence \(v\),
\[
 (1+|j|)^m\le(1+|i|)^m(1+d(i,j))^m
\]
shows that convolution by an exponentially summable
kernel is absolutely defined and preserves polynomial
growth. Exponential summability also permits the two
convolutions in the inverse identities to be interchanged.
Those identities therefore hold on polynomial-growth
sequences. In particular \(\mathcal A\) has no nonzero
polynomial-growth nullvector.

For a mixed kernel of any fixed finite degree, bound
each input polynomial by a polynomial in the output
coordinate times a polynomial in the joining-tree lengths.
The exponential tree bound absorbs the latter factors.
Thus the nonlinear Taylor coefficients are well defined.
At order \(n\) the stationary equation is
\[
 \mathcal A z_n=h_n(w_1,\ldots,w_n,z_1,\ldots,z_{n-1}),
\]
and has the unique polynomial-growth solution
\(\mathcal A^{-1}h_n\). These are also the equations
obtained by differentiating the analytic bounded-field
solution. Induction proves the assertion.
\end{proof}

For the actual constrained stationary branch at a fixed
endpoint, these are established inputs: the flat retained
precision has its exponential inverse, the affine
constraint chart is local, and the classical scalar
and stationary support families have exponentially
localized mixed derivatives. The finite fiber dimension
is per block; the whole lattice is not finite-dimensional
in this application. Only finite-order jets are being
extended.

For the bare Wilson action, the rooted affine candidate
above solves the constraint, contour and stationary
equations coefficientwise, with unconstrained current
and multiplier equal to zero. Lemma~\ref{cdmom-polynomial-lemma}
therefore identifies it with the actual stationary
section's polynomial jets. In particular this applies
through the two background orders needed below.

\subsection{The effective classical action is also stationary on this family}
\label{cdmom-effective-section}

The effective-action interface can be checked using the
actual envelope identity, without replacing its Hessian
by a bare Wilson Hessian.

Fix any permitted finite homogeneous depth \(M\).
The action \(S_M\) is the actual local minimum value of
the bare Wilson action under \(F_M\). Its differential
in left-relative output coordinates equals the target
multiplier of that minimum. This identity is an equality
of the established analytic kernels. It extends to
polynomial input jets by their exponential bounds.

Apply the preceding bare-action argument at the fixed
endpoint \(M\), with fine curvature equal to the prescribed
current curvature divided by \(M^2\). The explicit fine
candidate is Wilson stationary, so its multiplier jets
are zero. Consequently
\begin{equation}
          D_{\rm left}S_M(U^{\rm affine}(t))=0
                 \quad\hbox{as formal jets through order two}.
 \label{cdmom-effective-stationarity}
\end{equation}
The identity holds in every current-link component;
it is not just a derivative along the one-parameter family.
Indeed the envelope identifies the whole target multiplier,
and the unconstrained fine Euler derivative is zero.

At the next fixed endpoint \(K\), take the current affine
field with curvature divided by \(K^2\), and root it by
the actual direct-\(K\) based gauge. The classical local
gauge identity, also extended coefficientwise by exponential
bounds, preserves \eqref{cdmom-effective-stationarity}.
This candidate satisfies the actual \(S_M\) constrained
stationary equation with zero multiplier through the same
orders. Apply Lemma~\ref{cdmom-polynomial-lemma} to the
retained precision for \(S_M\). The candidate has exactly
the actual current stationary jets, and
\begin{equation}
             \lambda_{M,K}(t)=O(t^3)
 \label{cdmom-multiplier-jets}
\end{equation}
in this formal sense.

All assertions so far fix each endpoint when invoking
its inverse. They do not infer a gap uniform in \(M\)
or \(K\) for a one-shot fine operator. However, the
current-coordinate classical and stationary Taylor
kernels at fixed \(K\) have the already established
uniform bounds in preceding depth. Their polynomial
moments consequently pass to any varying-depth
structural limit by exponential domination.
Thus \eqref{cdmom-effective-stationarity} and
\eqref{cdmom-multiplier-jets} also hold in those joint
systems. The effective Hessian \(s''\) still depends
on \(S_M\); this result does not make it the bare
Wilson Hessian.

\subsection{An exact constant-curvature expression for the Maxwell coefficient}
\label{cdmom-determinant-section}

Fix a trace-orthonormal Cartan color \(H\), and distinct
directions \(\mu,\nu\). The output probe is
\[
 w_\mu(n;t)=t n_\nu H,\qquad w_\sigma=0\quad(\sigma\ne\mu).
\]
Its unrooted current affine representative is
\begin{equation}
 U_\mu^{\rm aff}(x;t)=\exp(i t x_\nu H/K^2),\qquad
 U_\sigma^{\rm aff}=I\quad(\sigma\ne\mu).
 \label{cdmom-probe}
\end{equation}
By \eqref{cdmom-average}, its output is exactly the
chosen formal probe. Its oriented \(\mu,\nu\) plaquette
has angle \(-tH/K^2\).

On this family write
\begin{align}
 G(t)&=D_b(U^{\rm aff}(t))^TD_b(U^{\rm aff}(t)),\notag\\
 V_a(t)&=s''(U^{\rm aff}(t))
           +D_b(U^{\rm aff}(t))D_b(U^{\rm aff}(t))^T
           +aC(U^{\rm aff}(t))^TC(U^{\rm aff}(t)),\notag\\
 P(t)&=q_{\rm eff}(w(t))+aI.
 \label{cdmom-probe-matrices}
\end{align}
All are formal matrix jets around their flat positive
operators. In particular no invertibility or positivity
on a nonzero global magnetic background is assumed.
Equation \eqref{cdmom-multiplier-jets} removes the
\(-\lambda\cdot f''\) term through second order.
The remaining Hessian \(s''\) retains the actual action.

Use the same block-zero site and bond assignments as
in the diagonal-density construction. Its pointwise
gauge invariance allows removal of the root gauge,
coefficientwise: the local gauge identities extend to
these polynomial gauge jets by the same exponential
summation. The density--moment lemma therefore gives
\begin{equation}
 \boxed{\begin{split}
 \beta(D_{M,K})=[t^2]\bigg\{
 &2\sum_{\substack{z\ {\rm noncenter}\\z\ {\rm in\ block}\ 0}}
                   \operatorname{tr}_{\mathfrak g}(\log G(t))_{zz}\\
 &-\sum_{\substack{e\ {\rm current}\\e_-\ {\rm in\ block}\ 0}}
                   \operatorname{tr}_{\mathfrak g}(\log V_a(t))_{ee}\\
 &+\sum_{\substack{c\ {\rm output}\\c_-=0}}
                   \operatorname{tr}_{\mathfrak g}(\log P(t))_{cc}
                      \bigg\}.
 \end{split}}
 \label{cdmom-exact-coefficient}
\end{equation}
Field-independent subtractions can be omitted inside
\([t^2]\). The factors \(2,-1,+1\) follow from
\(\beta=2[t^2]\rho_0\), and the determinant weights
\(1,-\tfrac12,+\tfrac12\).

Every coefficient on the right is defined by the
localized matrix-logarithm derivatives at zero.
Their source kernels absorb the polynomial weights,
so the sums are convergent at fixed endpoint.
For \(M=1\), \(s''\) in \(V_a(t)\) is the actual bare
Wilson Hessian on an explicit commuting constant-curvature
background. The coarse contribution is uniformly bounded
in \(M,K\) by Notebook 32. Hence the remaining ultraviolet
comparison can be made directly on the first two terms
of \eqref{cdmom-exact-coefficient}.

\subsection{What this settles and what it does not}
\label{cdmom-scope}

The formula matches an explicit constant-curvature trace
calculation to the actual Maxwell extraction, with the
actual composite average, stationary jets and measure
factors. It removes the multiplier-curvature jets on
this probe. It does not replace an arbitrary local
density by its value on an affine field: the individual
gauge identity of the chosen diagonal densities was
essential.

On a fixed periodic torus, constant nonzero Cartan flux
has quantization conditions. No analytic family of
periodic fluxes through zero is used here. Nor is an
unbounded affine field substituted into the original
bounded-field map. The identities are polynomial
extensions of localized derivatives in an infinite
joint system.

The fine vector--ghost expression still needs a bound
on its difference from a precisely normalized reference
logarithm, uniform as \(K\) grows. The fixed-endpoint
logarithm construction supplies no such bound and
computes no sign. Mixed and large-field histories,
the continuum observable construction and reconstruction
remain separate requirements for full existence.
